\documentclass[12pt,a4paper]{article}
\usepackage{geometry}
\geometry{margin=1.8cm}
\usepackage{fontspec}
\usepackage{unicode-math}
\usepackage{amsmath}
\usepackage{array}
\usepackage{booktabs}
\usepackage{enumitem}
\usepackage{graphicx}
\usepackage{xcolor}
\usepackage{fancyhdr}
\usepackage{titlesec}
\usepackage{setspace}
\setlength{\parindent}{0pt}
\setlength{\parskip}{0.6em}
\setlist[itemize]{leftmargin=1.5em, itemsep=0.4em}
\pagestyle{empty}
\setmainfont{Latin Modern Roman}
\setmathfont{Latin Modern Math}
\titleformat{\section}{\Large\bfseries}{}{0pt}{}

\newcommand{\ExamHeader}{%
\noindent
\begin{minipage}[t]{0.18\textwidth}
\fbox{\begin{minipage}[c][1.7cm][c]{2.4cm}\centering
LOGO\par
Dummy
\end{minipage}}
\end{minipage}
\hfill
\begin{minipage}[t]{0.78\textwidth}
\raggedright
\small
Fakülte: Mühendislik \\ Bölüm: ……………………………….……………… \\ Ders Kodu/İsmi: MAT 202 Diferansiyel Denklemler \\ Öğretim Elemanı: …………………………………………………
\par\medskip
\centering\bfseries
Sınav Evrakı \\ BÜTÜNLEME SINAVI \\ Tarih: 02/01/2026 \\ Süre: 90 dakika \\ Öğrenci No \\ Öğrenci Ad - Soyad \\ İmza
\end{minipage}
\par\medskip
\hrule
\vspace{0.8em}
}

\begin{document}
\ExamHeader
\textit{Sayfa 1/6}

\begin{center}
{\Large\bfseries SINAV KURALLARI}
\end{center}

\begin{itemize}
\item HESAP MAKİNESİ KULLANILAMAZ.
\item TELEFON KULLANILAMAZ.
\item AKILLI CİHAZ KULLANILAMAZ.
\item İŞLEMLERİNİZİ ANLAŞILIR ŞEKİLDE YAZINIZ.
\end{itemize}

\vfill
\begin{center}
{\large\bfseries SORULAR}
\end{center}
\clearpage
\ExamHeader
\pagestyle{empty}
\textit{Sayfa 2/6}

\textbf{Soru 1}

\[y^{2}\sin x\ \frac{dy}{\,dx}\  = \ {(y}^{3} + 2)\ \cos x\]

a) Yukarıdaki diferansiyel denklem türdeş (homojen) midir? Açıklayınız. (3 p)

b) Mertebesi (basamağı) nedir? (3 p)

c) Bu diferansiyel denklemi çözünüz (12 p).

d) Bulduğunuz çözüm bir çözüm ailesi midir, tek bir çözüm fonksiyonu mudur? Açıklayınız. (2 p)
\clearpage
\ExamHeader
\pagestyle{empty}
\textit{Sayfa 3/6}

\textbf{Soru 2}

\({x^{2}y}' + 2xy - y^{3} = \ 0 , \ x > 0\) , \(y(1) = 1\)

Aşağıdaki diferansiyel denklemin genel çözümünü Bernoulli diferansiyel denklemi (yöntemi) kullanarak bulunuz. Başka bir yöntem kullanılması durumunda puan verilmeyecektir.
\clearpage
\ExamHeader
\pagestyle{empty}
\textit{Sayfa 4/6}

\textbf{Soru 3}

Aşağıdaki başlangıç değer problemini çözünüz:

\({(x}^{2} + {\ 3y}^{2})dx - 2xy\,dy = 0,\) \(y(1) = 2\)
\clearpage
\ExamHeader
\pagestyle{empty}
\textit{Sayfa 5/6}

\textbf{Soru 4}

Aşağıdaki diferansiyel denklemi çözünüz:

\[y''' + y' = \sin x\]
\clearpage
\ExamHeader
\pagestyle{empty}
\textit{Sayfa 6/6}

\textbf{Soru 5}

Aşağıdaki diferansiyel denklemi çözünüz: (\(y^{(4)} = y^{(4)}\))

\[y^{(4)} - y''' + y'' - y' = x\]

\end{document}
